\documentclass[11pt,a4paper]{article}
\usepackage[margin=2.5cm]{geometry}
\usepackage{booktabs}
\usepackage{hyperref}
\usepackage{enumitem}
\usepackage{parskip}
\usepackage{microtype}

\setlength{\parskip}{0.55em}
\setlength{\parindent}{0pt}

\title{\textbf{Advanced NLP -- Exercise 2}\\[0.3em]
       \large Part 1}
\author{Netanel Azran}
\date{June 2026}

\begin{document}
\maketitle

% ============================================================
\section*{Question 1 \;---\; Societal Bias Beyond Gender}
% ============================================================

\textbf{Paper:} Moin Nadeem, Anna Bethke, and Siva Reddy.
``StereoSet: Measuring stereotypical bias in pretrained language models.''
\emph{ACL-IJCNLP 2021}.

StereoSet is a crowdsourced benchmark for measuring how much pretrained language
models rely on social stereotypes. It covers four bias categories --- race,
religion, profession, and gender --- with roughly 17,000 items in total.
Each item presents a short context followed by three candidate continuations:
a \emph{stereotypical} one, an \emph{anti-stereotypical} one, and a semantically
\emph{unrelated} one. For example, given ``\emph{Arabs are\ldots},'' the options
might be ``violent and dangerous,'' ``peaceful and welcoming,'' and ``made of
bricks.''

The paper proposes two metrics. The \emph{Stereotype Score} (SS) measures how
often the model prefers the stereotypical continuation over the
anti-stereotypical one (50 = unbiased). The \emph{Language Model Score} (LMS)
checks that the model still prefers meaningful text over nonsense, serving as a
sanity check. The two are combined into the \emph{ICAT score}, which rewards
models that are both capable and minimally biased. Tested models (BERT,
RoBERTa, XLNet, GPT-2) all show clear stereotypical preferences in the race and
religion subsets.

% ============================================================
\section*{Question 2 \;---\; Three Efficiency Methods}
% ============================================================

% ------ Method 1 ------
\subsection*{Method 1: Quantization \quad\textnormal{\small(Inference)}}

Quantization reduces the numerical precision of model weights --- from float32
down to float16, INT8, or INT4. This can be applied after training
(\emph{post-training quantization}) without any retraining, making it very
cheap to deploy.

The main saving is \emph{memory}: a 7B-parameter model shrinks from $\sim$28\,GB
in float32 to $\sim$3.5\,GB in INT4, which is the difference between needing a
data-centre GPU and fitting on a consumer card. Inference speed also improves
because lower-precision arithmetic is faster on modern hardware.

The trade-off is a small drop in accuracy, which becomes more noticeable at
very aggressive quantization (4-bit). Quantization-aware training recovers most
of the loss but requires an extra training phase.

% ------ Method 2 ------
\subsection*{Method 2: LoRA \quad\textnormal{\small(Training)}}

LoRA (Hu et al., 2022) fine-tunes a model by freezing all pretrained weights
and injecting small low-rank adapter matrices into each transformer layer. Only
those adapters are updated; the rest of the model is untouched. After training,
the adapters can be merged back into the base weights, so inference cost does
not change.

For a 7B model, the number of trainable parameters drops from 7 billion to
roughly 20--30 million (under 0.5\%), making it practical to fine-tune large
models on a single consumer GPU.

The trade-off is expressiveness: because weight updates are constrained to a
low-rank subspace, LoRA may fall slightly behind full fine-tuning on some tasks,
though in practice the gap is small.

% ------ Method 3 ------
\subsection*{Method 3: Knowledge Distillation \quad\textnormal{\small(Modeling)}}

Knowledge distillation (Hinton et al., 2015) trains a small \emph{student} model
to mimic a large \emph{teacher} model. Instead of training on hard one-hot
labels, the student matches the teacher's full output distribution (``soft''
probabilities), which carries richer information about similarity between
classes. A well-known NLP example is DistilBERT --- a 66M-parameter student
trained from BERT-base (110M parameters) that retains $\sim$97\% of BERT's
GLUE performance while being 40\% smaller and 60\% faster.

The trade-off is extra training cost: generating soft labels requires running
the large teacher, so distillation takes more compute than training the student
from scratch.

% ------ Pairwise ------
\subsection*{Pairwise Combinations}

\paragraph{Quantization + LoRA.}
These two are already combined in practice as \textbf{QLoRA} (Dettmers et al.,
2023). The frozen base model is quantized to 4-bit (saving the memory needed to
hold it), while the LoRA adapters are trained in 16-bit. This makes it possible
to fine-tune a 65B model on a single 48\,GB GPU --- something neither method
achieves alone.

\paragraph{Quantization + Knowledge Distillation.}
A natural pipeline: first distill the large teacher into a smaller student,
then quantize the student for deployment. The two savings stack --- a student
that is 4$\times$ smaller than the teacher, then quantized to INT8, ends up
using about one-eighth the memory of the original. The only cost is two
sequential training phases.

\paragraph{LoRA + Knowledge Distillation.}
These methods target different problems --- LoRA reduces fine-tuning cost while
distillation reduces model size --- so they are complementary rather than
redundant. One practical combination is to distill a large model into a smaller
student, then use LoRA to adapt that student cheaply to a downstream task.
Useful when both deployment efficiency and training efficiency are required at
the same time.

\end{document}
